% from main.tex
%%%%%%%%%%  scan idx 560  |  folio 547  %%%%%%%%%%
 i:\Pic(X)\longrightarrow K^\circ(X)
\end{equation}
l'application qui associe à la classe \(\cl_P(L)\in\Pic(X)\) d'un \(\mathcal O_X\)-module inversible \(L\) l'élément \(\cl^\circ(L)\) de \(K^\circ(X)\). Cette application est évidemment un homomorphisme de \(\Pic(X)\) dans le groupe multiplicatif \(K^\circ(X)^*\) des éléments inversibles de \(K^\circ(X)\). Par la suite nous noterons multiplicativement la loi de groupe de \(\Pic(X)\).

L'application \(i\) a un inverse à gauche défini par le foncteur «~det~» : Lorsque \(E\) est un \(\mathcal O_X\)-module localement libre de rang \(n\), on a
\[
 \det(E)=\bigwedge^n E.
\]
Si
\begin{equation}\tag{5.1.2}
0\longrightarrow E'\longrightarrow E\longrightarrow E''\longrightarrow0
\end{equation}
est une suite exacte de \(\mathcal O_X\)-modules localement libres, il y a un isomorphisme canonique
\[
 \det(E)\xrightarrow{\sim}(\det(E'))\otimes_{\mathcal O_X}(\det(E''))
\]
provenant de (6.2.9) et du fait que (5.1.2) est localement scindée. On obtient donc un homomorphisme de groupes abéliens
\begin{equation}\tag{5.1.3}
 \det:K^\circ(X)\longrightarrow\Pic(X)
\end{equation}
qui est l'inverse à gauche cherché de \(i\). En particulier \(i\) est injective, et nous identifierons par la suite \(\Pic(X)\) avec son image dans \(K^\circ(X)^*\).

\subsection*{5.2}
On définit une autre application
\begin{equation}\tag{5.2.1}
 \delta:\Pic(X)\longrightarrow\Filt^1(X)
\end{equation}
